\documentclass[../../main.tex]{subfiles}

\begin{document}

\section{函数}\label{section:ch01-函数}

函数这个概念我们曾多次遇到过，它在整个数学中占有重要的地位，提醒读者注意这一点很有必要。在这一节中我们给出它的精确的数学定义，并探讨与之有关的一些概念。

通常把函数视作某个法则，按照这个法则为集合 $A$ 的每一个元素确定集合 $B$ 的一个元素。在微积分中，函数常常用一个简单的公式给出，比如 $f(x) = 3x^{2} + 2$，有时式子比较复杂，如
\begin{align*}
f(x) = \sum_{k=1}^{\infty} x^{k}.
\end{align*}
人们常常并不明确地指出集合 $A$ 和 $B$ 是什么，而约定 $A$ 是使法则有意义的实数集，$B$ 就是整个实数集。

然而，随着数学的进展，人们需要对函数给出更为确切的定义。数学家们对于上述有关函数的提法尽管予以认同，但是，他们所使用的函数的定义更为精确。首先，给出以下定义。

\begin{definition}[指派法则]\label{definition:ch01-指派法则}
指派法则（rule of assignment）是两个集合的笛卡儿积 $C \times D$ 的一个子集 $r$，该子集满足这样的条件：$C$ 的每一个元素最多是 $r$ 中一个有序偶对的第一个坐标。
\end{definition}

这样，$C \times D$ 的一个子集 $r$ 如果满足
\begin{align*}
[(c, d) \in r \text{ 并且 } (c, d') \in r] \Longrightarrow [d = d'],
\end{align*}
则成为一个指派法则。我们可以将 $r$ 设想成一种指派方法，为 $C$ 的元素 $c$ 配置 $D$ 的满足 $(c, d) \in r$ 的元素 $d$。

对于一个指派法则 $r$，其定义域（domain）是由 $r$ 的元素的所有第一个坐标组成的 $C$ 的子集，其像集（image set）是由 $r$ 的元素的所有第二个坐标组成的 $D$ 的子集，即
\begin{align*}
r \text{ 的定义域} = \{c \mid \text{存在 } d \in D, \text{使得 } (c, d) \in r\},
\end{align*}
\begin{align*}
r \text{ 的像集} = \{d \mid \text{存在 } c \in C, \text{使得 } (c, d) \in r\}.
\end{align*}
如果给定了一个指派法则 $r$，其定义域和像是完全确定的。

现在来定义函数。

\begin{definition}[函数]\label{definition:ch01-函数}
函数（function）$f$ 是一个指派法则 $r$，连同一个包含 $r$ 的像集的集合 $B$。法则 $r$ 的定义域 $A$，称为函数 $f$ 的定义域（domain）。$r$ 的像集就称为 $f$ 的像集（image set）。集合 $B$ 称为 $f$ 的值域（range）。此外，对应、映射都是函数的同义语。
\end{definition}

若 $f$ 是以 $A$ 为定义域、$B$ 为值域的一个函数，我们就写作
\begin{align*}
f: A \longrightarrow B,
\end{align*}
读作“$f$ 是一个从 $A$ 到 $B$ 的函数”，或“$f$ 是从 $A$ 到 $B$ 中的一个映射”，或者简单地说“$f$ 将 $A$ 映到 $B$ 中”。有时把 $f$ 形象地看成是将 $A$ 中的点移到 $B$ 中去的几何变换。

如果 $f: A \longrightarrow B$，并且 $a$ 是 $A$ 的一个元素，那么 $f(a)$ 表示在法则 $f$ 下 $B$ 中给 $a$ 配置的那个唯一的元素，称为 $f$ 在 $a$ 点的值（value），有时也称为 $a$ 在 $f$ 下的像（image）或像集。正式地说，如果 $r$ 是函数 $f$ 的法则，$f(a)$ 就是使 $(a, f(a)) \in r$ 的 $B$ 的那个唯一的元素。

使用这种表示法，可以把前面几乎所有的函数表达得更加严谨。例如，我们可以说（用 $\mathbb{R}$ 表示实数集）：

\begin{quote}
“$f$ 是一个函数，它的法则是 $\{(x, x^{3} + 1) \mid x \in \mathbb{R}\}$，值域是 $\mathbb{R}$，”

也可以等价地说：

“$f: \mathbb{R} \longrightarrow \mathbb{R}$ 是一个函数，使得 $f(x) = x^{3} + 1$。”
\end{quote}

两种说法都相当精确地描述了同一个函数。但是，如果说成“$f$ 是函数，$f(x) = x^{3} + 1$”就很不妥当，因为它既没有指出 $f$ 的定义域，也没有指出它的值域。

\begin{definition}[限制]\label{definition:ch01-限制}
设 $f: A \longrightarrow B$，$A_{0}$ 为 $A$ 的一个子集，$f$ 在 $A_{0}$ 上的限制（restriction）定义为将 $A_{0}$ 映到 $B$ 中的一个函数，其法则是
\begin{align*}
\{(a, f(a)) \mid a \in A_{0}\}.
\end{align*}
记作 $f \mid A_{0}$，读成“$f$ 在 $A_{0}$ 上的限制”。
\end{definition}

\par
\begin{figure}[htbp]
\centering
\includegraphics[width=0.85\textwidth]{ch01-fig2-1}
\caption{同一公式在不同定义域与值域下的四个函数}\label{figure:ch01-四个平方函数}
\end{figure}

\begin{example}\label{example:ch01-函数-1}
设 $\mathbb{R}$ 为实数集，$\bar{\mathbb{R}}_{+}$ 为非负实数集。考虑函数
\begin{align*}
f: \mathbb{R} \longrightarrow \mathbb{R} \quad &\text{使得 } f(x) = x^{2},\\
g: \bar{\mathbb{R}}_{+} \longrightarrow \mathbb{R} \quad &\text{使得 } g(x) = x^{2},\\
h: \mathbb{R} \longrightarrow \bar{\mathbb{R}}_{+} \quad &\text{使得 } h(x) = x^{2},\\
k: \bar{\mathbb{R}}_{+} \longrightarrow \bar{\mathbb{R}}_{+} \quad &\text{使得 } k(x) = x^{2}.
\end{align*}
则函数 $g$ 与 $f$ 不同，因为它们的法则是 $\mathbb{R} \times \mathbb{R}$ 的不同子集。$g$ 是 $f$ 在 $\bar{\mathbb{R}}_{+}$ 上的限制。$h$ 与 $f$ 也不同，虽然它们的法则是同一个集合，但是值域不同。$k$ 与 $f$，$g$，$h$ 都不同。这些函数如\cref{figure:ch01-四个平方函数}所示。

\end{example}

限制函数的定义域和改变函数的值域，这是由已知函数做出新函数的两个方法。还有一种方法就是做出两个函数的复合函数。

\begin{definition}[复合]\label{definition:ch01-复合}
已知函数 $f: A \longrightarrow B$ 和 $g: B \longrightarrow C$，$f$ 与 $g$ 的复合（composition）$g \circ f$ 是一个函数 $g \circ f: A \longrightarrow C$，定义为 $(g \circ f)(a) = g(f(a))$。
\end{definition}

正式定义为 $g \circ f: A \longrightarrow C$ 是这样的一个函数，它的法则是
\begin{align*}
\{(a, c) \mid \text{对于某一个 } b \in B, \text{有 } f(a) = b \text{ 和 } g(b) = c\}.
\end{align*}
复合函数 $g \circ f$ 常常可以理解成从点 $a$ 到点 $f(a)$，然后再到 $g(f(a))$ 的一个物理运动，如\cref{figure:ch01-复合}所示。

\par
\begin{figure}[htbp]
\centering
\includegraphics[width=0.7\textwidth]{ch01-fig2-2}
\caption{函数的复合}\label{figure:ch01-复合}
\end{figure}

注意，$g \circ f$ 仅仅在 $f$ 的值域等于 $g$ 的定义域时才有意义。

\begin{example}\label{example:ch01-复合例}
由 $f(x) = 3x^{2} + 2$ 给出的函数 $f: \mathbb{R} \longrightarrow \mathbb{R}$ 与由 $g(x) = 5x$ 给出的函数 $g: \mathbb{R} \longrightarrow \mathbb{R}$ 的复合是函数 $g \circ f: \mathbb{R} \longrightarrow \mathbb{R}$，其定义为
\begin{align*}
(g \circ f)(x) = g(f(x)) = g(3x^{2} + 2) = 5(3x^{2} + 2).
\end{align*}
这时还可以作出复合 $f \circ g$。函数 $f \circ g: \mathbb{R} \longrightarrow \mathbb{R}$ 与 $g \circ f$ 完全不同，它定义为
\begin{align*}
(f \circ g)(x) = f(g(x)) = f(5x) = 3(5x)^{2} + 2.
\end{align*}
\end{example}

\begin{definition}[单射、满射、双射]\label{definition:ch01-单满双射}
函数 $f: A \longrightarrow B$ 称为单的（injective），如果 $A$ 中不同的点在 $f$ 下的像不相同。函数 $f: A \longrightarrow B$ 称为满的（surjective）（或 $f$ 将 $A$ 映满 $B$），如果 $B$ 的每一个元素都是 $A$ 中某元素在 $f$ 下的像。如果 $f$ 既是单的又是满的，则称 $f$ 为既单且满的（bijective），或一一的（one-to-one）。此外单的函数称为单射（injection），满的函数称为满射（surjection），一一的函数称为一一对应（one-to-one correspondence，one-to-one map）。
\end{definition}

更正式的定义是，如果
\begin{align*}
[f(a) = f(a')] \Longrightarrow [a = a'],
\end{align*}
则 $f$ 是单射。如果
\begin{align*}
[b \in B] \Longrightarrow [\text{对于至少一个 } a \in A, \text{使得 } b = f(a)],
\end{align*}
则 $f$ 是满射。

$f$ 是否是单射仅仅依赖于 $f$ 的法则，它是否是满射则还要依赖于 $f$ 的值域。不难验证：两个单射的复合是一个单射，两个满射的复合是一个满射。从而，两个一一对应的复合是一个一一对应。

如果 $f$ 是一个一一对应，则存在从 $B$ 到 $A$ 的一个函数，称为 $f$ 的逆（inverse），记作 $f^{-1}$，其元素 $f^{-1}(b)$ 定义为 $A$ 中满足 $f(a) = b$ 的那个唯一元素 $a$。对于 $b \in B$，由于 $f$ 是满射，可见这样的元素 $a \in A$ 是存在的；由于 $f$ 是单射，可见仅有一个这样的元素 $a \in A$。显然，如果 $f$ 是一一对应，则 $f^{-1}$ 也是一一对应。

\begin{example}\label{example:ch01-单满双射例}
仍考虑\cref{figure:ch01-四个平方函数}中的函数 $f$，$g$，$h$，$k$。函数 $f: \mathbb{R} \longrightarrow \mathbb{R}$，$f(x) = x^{2}$，既不是单射也不是满射。$f$ 在非负实数集上的限制 $g$ 是单射但不是满射。通过改变 $f$ 的值域而得到的函数 $h: \mathbb{R} \longrightarrow \bar{\mathbb{R}}_{+}$ 是满射但不是单射。通过限制 $f$ 的定义域并改变其值域而得到的函数 $k: \bar{\mathbb{R}}_{+} \longrightarrow \bar{\mathbb{R}}_{+}$ 既是单射又是满射，于是它有一个逆，通常称之为平方根函数。
\end{example}

为了判定一个函数是一一对应的，一个常用的方法是以下的引理，其证明留作习题。

\begin{lemma}[逆的存在性]\label{lemma:ch01-逆存在}
设 $f: A \longrightarrow B$。如果存在函数 $g: B \longrightarrow A$ 和 $h: B \longrightarrow A$，使得对于 $A$ 中每一个 $a$，$g(f(a)) = a$，并且对于 $B$ 中每一个 $b$，$f(h(b)) = b$，则 $f$ 是一个一一对应，并且 $g = h = f^{-1}$。
\end{lemma}

\begin{definition}[像与原像]\label{definition:ch01-像与原像}
设 $f: A \longrightarrow B$。$A_{0}$ 为 $A$ 的一个子集，用 $f(A_{0})$ 表示 $A_{0}$ 中的点在 $f$ 下的像的集合，这个集合称为 $A_{0}$ 在 $f$ 下的像（image）。正式的定义是
\begin{align}
f(A_{0}) = \{b \mid \text{对于某 } a \in A_{0}, b = f(a)\}.\label{align:ch01-像定义}
\end{align}
另一方面，若 $B_{0}$ 为 $B$ 的一个子集，用 $f^{-1}(B_{0})$ 表示 $A$ 中那些元素的集合，它们在 $f$ 下的像属于 $B_{0}$；$f^{-1}(B_{0})$ 称为 $B_{0}$ 在 $f$ 下的原像（preimage）（或 $B_{0}$ 的“反像”或“逆像”）。正式的定义是
\begin{align}
f^{-1}(B_{0}) = \{a \mid f(a) \in B_{0}\}.\label{align:ch01-原像定义}
\end{align}
\end{definition}

当然，有可能 $A$ 中任何点的像都不在 $B_{0}$ 中。这时，$f^{-1}(B_{0})$ 是空集。

注意，如果 $f: A \longrightarrow B$ 是一一对应，$B_{0} \subset B$，那么记号 $f^{-1}(B_{0})$ 就有两种含义：它既可以表示 $B_{0}$ 在 $f$ 下的原像，又可以表示 $B_{0}$ 在 $f^{-1}: B \longrightarrow A$ 下的像。但是这两种含义都恰好得出 $A$ 的同一个子集，所以不会造成混淆。

要正确使用记号 $f$ 与 $f^{-1}$，必须格外小心。例如，当 $f^{-1}$ 作用在 $B$ 的子集上时有很好的性质，它保持集合的“包含”、“并”、“交”与“差”。我们将时常利用这个事实。而 $f$ 作用在 $A$ 的子集上时，仅保持集合的“包含”和“并”。见习题 2 和 3。

另外一个值得注意的问题是，一般说来 $f^{-1}(f(A_{0})) = A_{0}$ 与 $f(f^{-1}(B_{0})) = B_{0}$ 并不成立（参见后面的例子）。请读者自行验证它们之间的下列关系：假定 $f: A \longrightarrow B$，$A_{0} \subset A$，和 $B_{0} \subset B$，则有
\begin{align*}
f^{-1}(f(A_{0})) \supset A_{0} \quad \text{和} \quad f(f^{-1}(B_{0})) \subset B_{0}.
\end{align*}
进一步，如果 $f$ 为单射，则第一个包含关系可改写为等式。如果 $f$ 为满射，则第二个包含关系可改写为等式。

\par
\begin{figure}[htbp]
\centering
\includegraphics[width=0.5\textwidth]{ch01-fig2-3}
\caption{$f(x) = 3x^{2} + 2$ 的图像}\label{figure:ch01-抛物线}
\end{figure}

\begin{example}\label{example:ch01-像原像例}
考虑由 $f(x) = 3x^{2} + 2$ 所定义的函数 $f: \mathbb{R} \longrightarrow \mathbb{R}$（见\cref{figure:ch01-抛物线}）。设 $[a, b]$ 为闭区间 $a \leqslant x \leqslant b$，则有
\begin{align*}
f^{-1}(f([0, 1])) = f^{-1}([2, 5]) = [-1, 1] \quad \text{和} \quad f(f^{-1}([0, 5])) = f([-1, 1]) = [2, 5].
\end{align*}

\end{example}

\subsection*{习题}

\begin{exercise}\label{exercise:ch01-函数-1}
设 $f: A \longrightarrow B$，$A_{0} \subset A$ 和 $B_{0} \subset B$。
\begin{enumerate}[(a)]
\item 证明 $A_{0} \subset f^{-1}(f(A_{0}))$，并且当 $f$ 是单射时，式中包含关系可换为等号。
\item 证明 $f(f^{-1}(B_{0})) \subset B_{0}$，并且当 $f$ 是满射时，式中包含关系可换为等号。
\end{enumerate}
\end{exercise}

\begin{exercise}\label{exercise:ch01-函数-2}
设 $f: A \longrightarrow B$，并且对于 $i = 0$ 和 $1$，$A_{i} \subset A$，$B_{i} \subset B$。证明 $f^{-1}$ 保持集合的包含、并、交和差：
\begin{enumerate}[(a)]
\item $B_{0} \subset B_{1} \Longrightarrow f^{-1}(B_{0}) \subset f^{-1}(B_{1})$.
\item $f^{-1}(B_{0} \cup B_{1}) = f^{-1}(B_{0}) \cup f^{-1}(B_{1})$.
\item $f^{-1}(B_{0} \cap B_{1}) = f^{-1}(B_{0}) \cap f^{-1}(B_{1})$.
\item $f^{-1}(B_{0} - B_{1}) = f^{-1}(B_{0}) - f^{-1}(B_{1})$.
\end{enumerate}
证明 $f$ 仅保持集合的包含与并：
\begin{enumerate}[(e)]
\item $A_{0} \subset A_{1} \Longrightarrow f(A_{0}) \subset f(A_{1})$.
\item $f(A_{0} \cup A_{1}) = f(A_{0}) \cup f(A_{1})$.
\item $f(A_{0} \cap A_{1}) \subset f(A_{0}) \cap f(A_{1})$. 证明当 $f$ 是单射时，式中包含关系 $\subset$ 可换为等号。
\item $f(A_{0} - A_{1}) \supset f(A_{0}) - f(A_{1})$，证明当 $f$ 是单射时，式中包含关系 $\supset$ 可换为等号。
\end{enumerate}
\end{exercise}

\begin{exercise}\label{exercise:ch01-函数-3}
证明习题\cref{exercise:ch01-函数-2}中的 (b)、(c)、(f)、(g) 小题对于任意并和任意交成立。
\end{exercise}

\begin{exercise}\label{exercise:ch01-函数-4}
设 $f: A \longrightarrow B$，$g: B \longrightarrow C$。
\begin{enumerate}[(a)]
\item 若 $C_{0} \subset C$，证明 $(g \circ f)^{-1}(C_{0}) = f^{-1}(g^{-1}(C_{0}))$.
\item 若 $f$ 和 $g$ 都是单射，证明 $g \circ f$ 也是一个单射。
\item 若 $g \circ f$ 是一个单射，试讨论 $f$ 和 $g$ 是否是单射？
\item 若 $f$ 和 $g$ 都是满射，证明 $g \circ f$ 也是一个满射。
\item 若 $g \circ f$ 是一个满射，那么试讨论 $f$ 和 $g$ 是否是满射？
\item 将从 (b)$\sim$(e) 各小题得到的结论总结为一个定理。
\end{enumerate}
\end{exercise}

\begin{exercise}\label{exercise:ch01-函数-5}
集合 $C$ 的恒等函数（identity function）一般记作 $i_{C}$。函数 $i_{C}: C \longrightarrow C$ 定义为：对于任意 $x \in C$，$i_{C}(x) = x$。对于函数 $f: A \longrightarrow B$，如果存在函数 $g: B \longrightarrow A$，使 $g \circ f = i_{A}$，则称 $g$ 为 $f$ 的左逆（left inverse）。如果存在函数 $h: B \longrightarrow A$，使 $f \circ h = i_{B}$，则称 $h$ 为 $f$ 的右逆（right inverse）。
\begin{enumerate}[(a)]
\item 证明，若 $f$ 有左逆，则 $f$ 是单射；若 $f$ 有右逆，则 $f$ 是满射。
\item 举出一个有左逆而无右逆的函数的例子。
\item 举出一个有右逆而无左逆的函数的例子。
\item 是否存在有多个左逆的函数？是否存在有多个右逆的函数？
\item 证明：若 $f$ 既有左逆 $g$，又有右逆 $h$，则 $f$ 是一一对应并且
\begin{align*}
g = h = f^{-1}.
\end{align*}
\end{enumerate}
\end{exercise}

\begin{exercise}\label{exercise:ch01-函数-6}
假设函数 $f: \mathbb{R} \longrightarrow \mathbb{R}$ 是通过式子 $f(x) = x^{3} - x$ 来定义的。适当限制其定义域及改变值域得到一个一一的函数 $g$。画出 $g$ 与 $g^{-1}$ 的图形。（$g$ 有多种可能的选择。）
\end{exercise}

\end{document}